\section{Congruence normal forms of symmetric and skew-symmetric unitaries}

This section treats unitaries $x$ related to their entrywise conjugates and
transposes. It proves the congruence normal forms $x=S^T\Lambda S$ of
symmetric and skew-symmetric unitaries
\cite[Lemma~\texttt{lemma:conjclass-normalform-continuous}]{Cirac2017MPU},
the real normal form of a skew-symmetric orthogonal matrix, and the gauges
that make a pair $u,v$ with $\overline u=(x\otimes y)u$ and
$\overline v=v(y^\dagger\otimes x^\dagger)$ real or symplectic.

\begin{lemma}[Scalar phase from a unitary relation]
    \label{thm:mpug_unitary_entrywise_scalar_phase}
    \lean{Matrix.scalar_mul_star_eq_one_of_mul_map_star_eq_smul_one}
    \leanok
    Let $I$ be a nonempty finite set, let $T,S\in\operatorname{Mat}_I(\C)$ be
    unitary, and let $\sigma\in\C$. As in
    \cite[lines~1557--1562]{FrancoRubio2025MPUGauging}, if the product of
    $T$ and the entrywise conjugate $S^*$ of $S$ satisfies
    \begin{align}
        TS^* & =\sigma\mathbf 1,
        \label{eq:mpug_unitary_entrywise_scalar_relation}
    \end{align}
    then the scalar $\sigma$ has unit modulus:
    \begin{align}
        \sigma\overline\sigma & =1.
        \label{eq:mpug_unitary_entrywise_scalar_phase}
    \end{align}
\end{lemma}

\begin{proof}\leanok
    The matrix $S^*$ is unitary, so $TS^*=\sigma\mathbf 1$ is unitary. Hence
    \begin{align}
        (\sigma\mathbf 1)(\sigma\mathbf 1)^\dagger
         & =\sigma\overline\sigma\,\mathbf 1=\mathbf 1.
    \end{align}
    Evaluating a diagonal entry gives
    (\ref{eq:mpug_unitary_entrywise_scalar_phase}).
\end{proof}

\begin{lemma}[Reciprocal scalars in the paired relations]
    \label{thm:mpug_unitary_entrywise_paired_scalars}
    \lean{Matrix.paired_scalars_mul_eq_one_of_mul_map_star_eq_smul_one}
    \leanok
    \uses{thm:mpug_unitary_entrywise_scalar_phase}
    Let $I$ be a nonempty finite set, let $T,S\in\operatorname{Mat}_I(\C)$ be
    unitary, and let $\sigma,\tau\in\C$. As in the relation
    $\sigma_g\sigma_{g^{-1}}=1$ of
    \cite[lines~1559--1563]{FrancoRubio2025MPUGauging}, if
    \begin{align}
        TS^* & =\sigma\mathbf 1,
        \label{eq:mpug_unitary_entrywise_rel_ts} \\
        ST^* & =\tau\mathbf 1,
        \label{eq:mpug_unitary_entrywise_rel_st}
    \end{align}
    then
    \begin{align}
        \sigma\tau & =1.
        \label{eq:mpug_unitary_entrywise_paired_scalars}
    \end{align}
\end{lemma}

\begin{proof}\leanok
    \uses{thm:mpug_unitary_entrywise_scalar_phase}
    Multiplying~(\ref{eq:mpug_unitary_entrywise_rel_ts}) by $T^{-1}$ on the
    left and $T$ on the right gives $S^*T=\sigma\mathbf 1$. Entrywise
    conjugation of~(\ref{eq:mpug_unitary_entrywise_rel_st}) gives
    $S^*T=\overline\tau\,\mathbf 1$. A diagonal entry therefore gives
    $\sigma=\overline\tau$. Applying
    Lemma~\ref{thm:mpug_unitary_entrywise_scalar_phase}
    to~(\ref{eq:mpug_unitary_entrywise_rel_st}) yields $\tau\overline\tau=1$,
    and hence $\sigma\tau=1$.
\end{proof}

\begin{lemma}[Sign dichotomy in the self relation]
    \label{thm:mpug_unitary_entrywise_self_sign}
    \lean{Matrix.scalar_eq_one_or_neg_one_of_mul_map_star_self_eq_smul_one}
    \leanok
    \uses{thm:mpug_unitary_entrywise_paired_scalars}
    Let $I$ be a nonempty finite set, let $T\in\operatorname{Mat}_I(\C)$ be
    unitary, and let $\sigma\in\C$. As in the sign step
    $\Id=xx^\dagger=e^{2i\phi}\Id$ of \cite[lines~1036--1038]{Cirac2017MPU} and
    the involution conclusion of
    \cite[lines~1563--1567]{FrancoRubio2025MPUGauging}, if
    \begin{align}
        TT^* & =\sigma\mathbf 1,
        \label{eq:mpug_unitary_entrywise_self_relation}
    \end{align}
    then
    \begin{align}
        \sigma & \in\{1,-1\}.
        \label{eq:mpug_unitary_entrywise_self_sign}
    \end{align}
\end{lemma}

\begin{proof}\leanok
    \uses{thm:mpug_unitary_entrywise_paired_scalars}
    Applying Lemma~\ref{thm:mpug_unitary_entrywise_paired_scalars} with $S=T$
    and $\tau=\sigma$ gives $\sigma^2=1$, so
    $\sigma=1$ or $\sigma=-1$.
\end{proof}

\begin{lemma}[Common transpose sign of two conjugation factors]
    \label{thm:mpu_conjugation_common_transpose_sign}
    \lean{Matrix.transpose_mul_map_star_of_mem_unitaryGroup,
        Matrix.eq_smul_transpose_of_map_star_mul_eq_smul_one,
        Matrix.exists_common_transpose_sign_of_kronecker_map_star_mul_eq_one,
        Matrix.transpose_eq_self_or_transpose_eq_neg_of_kronecker_map_star_mul_eq_one}
    \leanok
    \uses{thm:mpug_kronecker_identity_factors, thm:mpug_unitary_entrywise_self_sign}
    Let $x$ and $y$ be unitary matrices over $\mathbb C$, indexed by nonempty
    finite sets, whose entrywise conjugates satisfy
    \begin{align}
        \overline xx\otimes\overline yy & =\Id.
    \end{align}
    Then there is a sign $\sigma\in\{1,-1\}$ with
    \begin{align}
        x & =\sigma x^T, & y & =\sigma y^T.
    \end{align}
    Equivalently, either $x=x^T$ and $y=y^T$, or $x=-x^T$ and $y=-y^T$.
    This is the algebraic conclusion of \cite[lines~1033--1038]{Cirac2017MPU}.
\end{lemma}

\begin{proof}\leanok
    \uses{thm:mpug_kronecker_identity_factors, thm:mpug_unitary_entrywise_self_sign}
    A Kronecker product equal to the identity has reciprocal scalar factors, so
    $\overline xx=e^{i\phi}\Id$ and $\overline yy=e^{-i\phi}\Id$. Since $x$ is
    unitary, $x^T\overline x=\Id$; multiplying $\overline xx=e^{i\phi}\Id$ on
    the left by $x^T$ gives $x=e^{i\phi}x^T$, and similarly
    $y=e^{-i\phi}y^T$. Transposing $x=e^{i\phi}x^T$ gives $x^T=e^{i\phi}x$,
    hence $x=e^{2i\phi}x$; since $x$ is invertible, $e^{2i\phi}=1$, so
    $e^{i\phi}=\pm1$ and $e^{-i\phi}=e^{i\phi}$. The two matrices therefore
    carry one transpose sign.
\end{proof}

\begin{lemma}[Symmetric branch of the unitary congruence normal form]
    \label{lem:mpu_symmetric_unitary_congruence}
    \lean{Matrix.exists_symmetric_unitary_squareRoot}
    \leanok
    Let $x$ be a symmetric unitary matrix. There is a symmetric unitary matrix
    $S$ such that
    \begin{align}
        x & = S^T S.
    \end{align}
    This is the symmetric branch of
    \cite[Lemma~\texttt{lemma:conjclass-normalform-continuous}]{Cirac2017MPU}.
    % Source lines: paper_v2.tex 1054--1064.
\end{lemma}

\begin{proof}\leanok
    The finite spectral decomposition of $x$ takes the form
    \begin{align}
        x & = \sum_{\lambda\in\sigma(x)} \lambda P_\lambda;
    \end{align}
    set
    \begin{align}
        S & = \sum_{\lambda\in\sigma(x)} \sqrt{\lambda}\,P_\lambda.
    \end{align}
    Each spectral value has modulus one, so its square root has modulus one,
    and hence $S$ is unitary and $S^2=x$. On the finite spectrum, interpolate
    the square-root function by a polynomial $q$. Then $S=q(x)$, and
    $x^T=x$ implies $S^T=q(x^T)=q(x)=S$. Therefore $x=S^2=S^T S$.
\end{proof}

\begin{lemma}[Determinant of $i$ times a traceless Hermitian involution]
    \label{lem:mpu_det_i_hermitian_involution}
    \lean{Matrix.IsHermitian.det_I_smul_eq_one}
    \leanok
    Let $I$ be a finite set and let $H\in\operatorname{Mat}_I(\mathbb C)$ be
    Hermitian with $H^2=\Id$ and $\operatorname{tr}H=0$. Then
    \begin{align}
        \det(iH) & =1.
    \end{align}
\end{lemma}

\begin{proof}\leanok
    Let $\mu_1,\dots,\mu_{|I|}$ be the real eigenvalues of $H$. Applying
    $H^2=\Id$ to an eigenvector gives $\mu_j^2=1$, so $\mu_j=\pm1$ and
    $i\mu_j=e^{i\pi\mu_j/2}$. The trace condition gives $\sum_j\mu_j=0$.
    Hence
    \begin{align}
        \det(iH)
         & =\prod_j i\mu_j
        =\exp\Bigl(\frac{i\pi}{2}\sum_j\mu_j\Bigr)
        =1.
    \end{align}
\end{proof}

\begin{lemma}[Skew congruence from paired spectral projectors]
    \label{lem:mpu_skew_unitary_paired_projectors}
    \lean{Matrix.exists_skew_unitary_congruence_of_paired_projectors}
    \leanok
    \uses{lem:mpu_det_i_hermitian_involution}
    Let $K$ and $I$ be finite sets, let $E_k\in\mathbb R$ for $k\in K$, and
    let $P_k$ be Hermitian projectors on $\mathbb C^I$. Suppose that the family
    $\{P_k,P_k^T\}_{k\in K}$ consists of mutually orthogonal projectors and
    sums to the identity. Define
    \begin{align}
        x & = \sum_{k\in K} e^{-iE_k}(P_k-P_k^T).
    \end{align}
    Then the matrices
    \begin{align}
        S & = e^{-i\pi/4}\sum_{k\in K} e^{-iE_k/2}(P_k+P_k^T),
    \end{align}
    and
    \begin{align}
        \widetilde\Lambda & = i\sum_{k\in K}(P_k-P_k^T)
    \end{align}
    are unitary, $S^T=S$, $\overline{\widetilde\Lambda}=\widetilde\Lambda$,
    $\widetilde\Lambda^T=-\widetilde\Lambda$, $\det\widetilde\Lambda=1$, and
    \begin{align}
        x & = S^T\widetilde\Lambda S.
    \end{align}
    This is the paired-projector construction for the skew branch of
    \cite[Lemma~\texttt{lemma:conjclass-normalform-continuous}]{Cirac2017MPU}.
    % Source lines: paper_v2.tex 1066--1083.
\end{lemma}

\begin{proof}\leanok
    \uses{lem:mpu_det_i_hermitian_involution}
    Viewing the two projectors in each pair as separate members of one
    complete orthogonal family shows that a weighted sum of this family is
    unitary whenever every weight has modulus one. This applies to $S$, whose
    two weights at $k$ both equal $e^{-i\pi/4}e^{-iE_k/2}$, and to
    $\widetilde\Lambda$, whose weights are $i$ and $-i$.
    Transposition exchanges the two members of each pair, proving the symmetry
    of $S$ and the skew symmetry of $\widetilde\Lambda$. Hermiticity gives
    $\overline{P_k}=P_k^T$, hence
    $\overline{\widetilde\Lambda}=\widetilde\Lambda$. Finally, orthogonality
    reduces the product to its weights, and
    \begin{align}
        \left(e^{-i\pi/4}e^{-iE_k/2}\right)^2 i & = e^{-iE_k}.
    \end{align}
    The transposed member has the negative weight, which gives
    $x=S^T\widetilde\Lambda S$. For the determinant, the same reduction gives
    $\widetilde\Lambda^2=-\Id$, so $\widetilde\Lambda=iH$ with
    $H=\sum_{k\in K}(P_k-P_k^T)$ a Hermitian involution. Since
    $\operatorname{tr}P_k^T=\operatorname{tr}P_k$, the trace of $H$ vanishes,
    and Lemma~\ref{lem:mpu_det_i_hermitian_involution} gives
    $\det\widetilde\Lambda=1$. These are the eigenvalue pairs $\pm i$ of the
    source.
\end{proof}

\begin{lemma}[Transpose pairing of spectral projectors]
    \label{lem:mpu_transpose_spectral_projector}
    \lean{Matrix.spectralProjector, Matrix.conjTranspose_spectralProjector,
        Matrix.spectralProjector_mul, Matrix.sum_spectralProjector,
        Matrix.sum_smul_spectralProjector,
        Matrix.transpose_spectralProjector_of_transpose_eq_neg}
    \leanok
    Let $x$ be a normal complex matrix with finite spectrum $\sigma(x)$, and for
    $\mu\in\mathbb C$ let $P_\mu$ be the spectral projector of $x$ at $\mu$,
    that is, the indicator function of the point $\mu$ applied to $x$. Then
    each $P_\mu$ is Hermitian,
    \begin{align}
        P_\mu P_\nu                     & =\delta_{\mu\nu}P_\mu, \\
        \sum_{\mu\in\sigma(x)}P_\mu     & =\Id,                  \\
        \sum_{\mu\in\sigma(x)}\mu P_\mu & =x.
    \end{align}
    If in addition $x^T=-x$, then
    \begin{align}
        P_\mu^T & =P_{-\mu}.
    \end{align}
    This relation pairs $P_E$ with $P_E^T$ in the skew branch of
    \cite[Lemma~\texttt{lemma:conjclass-normalform-continuous}]{Cirac2017MPU}.
    % Source lines: paper_v2.tex 1066--1070.
\end{lemma}

\begin{proof}\leanok
    The first three identities hold for the indicator functions on
    $\sigma(x)$. For the last, interpolate the indicator of $\mu$ on the
    finite set $\sigma(x)$ by a polynomial $q$, so that $P_\mu=q(x)$. Then
    $P_\mu^T=q(x^T)=q(-x)$. Since $\sigma(x)=\sigma(x^T)=-\sigma(x)$, the
    polynomial $z\mapsto q(-z)$ agrees on $\sigma(x)$ with the indicator of
    $-\mu$, and hence $q(-x)=P_{-\mu}$.
\end{proof}

\begin{lemma}[Paired spectral data of a skew-symmetric unitary]
    \label{lem:mpu_skew_unitary_paired_spectral_data}
    \lean{Matrix.exists_paired_spectralProjectors_of_transpose_eq_neg,
        Complex.arg_nonpos_iff_not_arg_neg_nonpos}
    \leanok
    \uses{lem:mpu_transpose_spectral_projector}
    Let $x$ be a unitary matrix with $x^T=-x$. There are a finite set $K$,
    Hermitian projectors $P_k$ and numbers $E_k\in[0,\pi)$, for $k\in K$,
    such that the family $\{P_k,P_k^T\}_{k\in K}$ consists of mutually
    orthogonal projectors summing to the identity, and
    \begin{align}
        x & = \sum_{k\in K} e^{-iE_k}(P_k-P_k^T).
    \end{align}
    This is the spectral decomposition in the skew branch of
    \cite[Lemma~\texttt{lemma:conjclass-normalform-continuous}]{Cirac2017MPU}.
    % Source lines: paper_v2.tex 1066--1070.
\end{lemma}

\begin{proof}\leanok
    \uses{lem:mpu_transpose_spectral_projector}
    Every spectral value of $x$ has modulus one, so it is nonzero and differs
    from its negative. By
    Lemma~\ref{lem:mpu_transpose_spectral_projector}, $\sigma(x)=-\sigma(x)$
    and $P_\lambda^T=P_{-\lambda}$. A nonzero complex number has argument in
    $(-\pi,0]$ exactly when its negative does not, so let $K$ be the set of
    spectral values with argument in $(-\pi,0]$; it contains exactly one
    member of each pair $\{\lambda,-\lambda\}$. For $\lambda\in K$ put
    $P_\lambda$ and $E_\lambda=-\arg\lambda\in[0,\pi)$. The spectral
    decomposition of $x$ then splits into the two members of each pair, with
    $-\lambda=-e^{-iE_\lambda}$ carrying the projector $P_\lambda^T$.
\end{proof}

\begin{lemma}[Skew branch of the unitary congruence normal form]
    \label{lem:mpu_skew_unitary_congruence}
    \lean{Matrix.exists_skew_unitary_congruence}
    \leanok
    \uses{lem:mpu_skew_unitary_paired_spectral_data,
        lem:mpu_skew_unitary_paired_projectors}
    Let $x$ be a unitary matrix with $x^T=-x$. There are a symmetric unitary
    matrix $S$ and a real skew-symmetric unitary matrix $\widetilde\Lambda$
    with $\det\widetilde\Lambda=1$ such that
    \begin{align}
        x & = S^T\widetilde\Lambda S.
    \end{align}
    This is the skew branch of
    \cite[Lemma~\texttt{lemma:conjclass-normalform-continuous}]{Cirac2017MPU}.
    % Source lines: paper_v2.tex 1066--1083.
\end{lemma}

\begin{proof}\leanok
    \uses{lem:mpu_skew_unitary_paired_spectral_data,
        lem:mpu_skew_unitary_paired_projectors}
    Lemma~\ref{lem:mpu_skew_unitary_paired_spectral_data} writes $x$ in the
    form required by Lemma~\ref{lem:mpu_skew_unitary_paired_projectors}, which
    gives $S$ and $\widetilde\Lambda$.
\end{proof}

\begin{lemma}[Normal form of symmetric and skew-symmetric unitaries]
    \label{lemma:conjclass-normalform-continuous}
    \lean{Matrix.exists_unitary_congruence_normalForm}
    \leanok
    \uses{lem:mpu_symmetric_unitary_congruence, lem:mpu_skew_unitary_congruence}
    Let $x$ be unitary and satisfy $x=x^T$ or $x=-x^T$. There are a symmetric
    unitary $S$ and a matrix $\widetilde\Lambda$ such that
    \begin{align}
        x & =S^T\widetilde\Lambda S, & S & =S^T.
    \end{align}
    In the symmetric case, $\widetilde\Lambda$ is real and symmetric. In the
    skew-symmetric case, it is real and skew-symmetric. In both cases,
    $\det(\widetilde\Lambda)=1$. This is
    \cite[Lemma~\texttt{lemma:conjclass-normalform-continuous}]{Cirac2017MPU}.
    % Source lines: paper_v2.tex 1041--1084.
\end{lemma}

\begin{proof}\leanok
    \uses{lem:mpu_symmetric_unitary_congruence, lem:mpu_skew_unitary_congruence}
    In the symmetric case, Lemma~\ref{lem:mpu_symmetric_unitary_congruence}
    gives $x=S^TS$, so $\widetilde\Lambda=\Id$ works. The skew-symmetric case
    is Lemma~\ref{lem:mpu_skew_unitary_congruence}. A matrix that is both
    symmetric and skew-symmetric vanishes, and a vanishing unitary exists
    only on the zero space, where both conditions on $\widetilde\Lambda$ hold.
\end{proof}

\begin{lemma}[Invariant orthogonal planes]
    \label{lem:mpu_skew_orthogonal_plane}
    \lean{LinearIsometry.orthonormal_pair_of_sq_eq_neg,
        LinearIsometry.apply_pair_eq_standardSkew_sum,
        LinearIsometry.mapsTo_span_pair_of_sq_eq_neg,
        LinearIsometry.mapsTo_orthogonal_of_sq_eq_neg,
        LinearIsometry.restrictInvariantOrthogonal_sq,
        Matrix.orthogonalLinearIsometry_sq}
    \leanok
    Let $J$ be a linear isometry of a real inner-product space satisfying
    $J^2=-\Id$. For every unit vector $v$, the ordered pair $(v,-Jv)$ is
    orthonormal. Its span and its orthogonal complement are invariant under
    $J$, and the restriction to the complement also squares to $-\Id$.
    On that ordered plane the matrix of $J$ is
    $Y=\left(\begin{smallmatrix}0&1\\-1&0\end{smallmatrix}\right)$.
    Every real skew-symmetric orthogonal matrix supplies such an isometry.
\end{lemma}

\begin{proof}\leanok
    Isometry and $J^2=-\Id$ give
    $\langle x,Jy\rangle=-\langle Jx,y\rangle$. Taking $x=y=v$ shows
    $\langle v,Jv\rangle=0$. The identities $Jv=-(-Jv)$ and
    $J(-Jv)=v$ give invariance and the displayed orientation. For an invariant
    subspace $W$, the same skew-adjoint identity gives
    $\langle w,Jz\rangle=-\langle Jw,z\rangle=0$ for $w\in W$ and
    $z\in W^\perp$, proving invariance of the complement. Finally,
    $A^T=-A$ and $AA^T=\Id$ imply $A^2=-\Id$.
\end{proof}

\begin{lemma}[Sign and indexing of the standard skew form]
    \label{lem:mpu_standard_skew_indexing}
    \lean{Matrix.interleavedSkewForm_eq_neg_J_submatrix,
        Matrix.interleavedSkewFormFin_apply}
    \leanok
    For a finite index set $I$, identify $(i,0),(i,1)\in I\times\{0,1\}$
    with $i$ in the first and second copies of $I\sqcup I$. Under this
    identification $\Id_I\otimes Y$ corresponds to
    $\left(\begin{smallmatrix}0&\Id_I\\-\Id_I&0\end{smallmatrix}\right)$,
    the negative of the split matrix
    $\left(\begin{smallmatrix}0&-\Id_I\\\Id_I&0\end{smallmatrix}\right)$.
    For $I=\{0,\ldots,m-1\}$ the interleaved enumeration gives the same
    ordered blocks on $\{0,\ldots,2m-1\}$.
\end{lemma}

\begin{proof}\leanok
    Compare the four choices of the two binary coordinates. The two diagonal
    blocks vanish, and the off-diagonal blocks are $\Id_I$ and $-\Id_I$
    in that order. Reindexing preserves these entries.
\end{proof}

\begin{lemma}[Paired orthonormal basis]
    \label{lem:mpu_skew_paired_orthonormal_basis}
    \lean{LinearIsometry.exists_paired_orthonormalBasis}
    \leanok
    \uses{lem:mpu_skew_orthogonal_plane}
    Let $E$ be a real inner-product space of dimension $2m$, and let $J$ be
    a linear isometry with $J^2=-\Id$. There is an orthonormal basis
    $(e_{i,0},e_{i,1})_{i=0}^{m-1}$ such that
    $Je_{i,0}=-e_{i,1}$ and $Je_{i,1}=e_{i,0}$.
\end{lemma}

\begin{proof}\leanok
    \uses{lem:mpu_skew_orthogonal_plane}
    Induct on $m$. The empty basis covers $m=0$. Otherwise choose a unit
    vector $v$ and use the ordered orthonormal pair $(v,-Jv)$.
    Its invariant orthogonal complement has dimension $2m-2$.
    Append its inductively constructed paired basis to this pair.
    The resulting orthonormal family has $2m$ members and therefore spans $E$.
\end{proof}

\begin{lemma}[Real skew-orthogonal normal form]
    \label{thm:mpu_real_skew_orthogonal_normal_form}
    \lean{Matrix.toMatrix_paired_orthonormalBasis,
        Matrix.exists_orthogonal_congruence_interleavedSkewForm}
    \leanok
    \uses{lem:mpu_skew_paired_orthonormal_basis, lem:mpu_standard_skew_indexing}
    Let $A\in\operatorname{Mat}_{2m}(\mathbb R)$ be orthogonal with
    $A^T=-A$. Set $\Sigma_m=\Id_m\otimes Y$, in interleaved order, where
    $Y=\left(\begin{smallmatrix}0&1\\-1&0\end{smallmatrix}\right)$.
    There is a real orthogonal matrix $O$ such that
    $O^T\Sigma_m O=A$.
    This is the real normal form used in the proof of
    \cite[Proposition~V.2]{Cirac2017MPU}.
    % Source lines: paper_v2.tex 1129--1140.
\end{lemma}

\begin{proof}\leanok
    \uses{lem:mpu_skew_paired_orthonormal_basis, lem:mpu_standard_skew_indexing}
    The action of $A$ is an isometry squaring to $-\Id$. In the paired
    orthonormal basis of Lemma~\ref{lem:mpu_skew_paired_orthonormal_basis},
    its matrix is $\Sigma_m$. The orthogonal change of coordinates from
    the original basis to this paired basis gives the displayed congruence.
\end{proof}

\begin{lemma}[Complex skew-unitary normal form]
    \label{thm:mpu_complex_skew_unitary_normal_form}
    \lean{Matrix.exists_unitary_congruence_interleavedSkewForm}
    \leanok
    \uses{lem:mpu_skew_unitary_congruence, thm:mpu_real_skew_orthogonal_normal_form}
    Let $x\in\operatorname{Mat}_{2m}(\mathbb C)$ be unitary with $x^T=-x$.
    There is a unitary matrix $U$ such that $x=U^T\Sigma_m U$.
    This is the complex normal form used in the proof of
    \cite[Proposition~V.2]{Cirac2017MPU}.
    % Source lines: paper_v2.tex 1129--1140.
\end{lemma}

\begin{proof}\leanok
    \uses{lem:mpu_skew_unitary_congruence, thm:mpu_real_skew_orthogonal_normal_form}
    The skew congruence factorization gives $x=S^T\Lambda S$ with $S$
    unitary and $\Lambda$ real, skew-symmetric and unitary.
    The real normal form gives $\Lambda=O^T\Sigma_m O$ with $O$ orthogonal.
    Regard $O$ as a complex unitary matrix and set $U=OS$.
\end{proof}

\begin{lemma}[Orthogonal and symplectic gauges of conjugation-covariant unitaries]
    \label{lem:mpu_conjugation_gate_gauge}
    \lean{Matrix.map_star_gaugeLeft_eq_of_congruence,
        Matrix.map_star_gaugeRight_eq_of_congruence, Matrix.sourceSkewForm,
        Matrix.exists_unitary_congruence_sourceSkewForm,
        Matrix.exists_orthogonal_gauge_of_symmetric,
        Matrix.exists_symplectic_gauge_of_skew,
        Matrix.exists_phase_of_kronecker_mul_map_star_eq,
        Matrix.exists_orthogonal_factors_of_kronecker_mul,
        Matrix.exists_symplectic_factors_of_kronecker_mul}
    \leanok
    \uses{lemma:conjclass-normalform-continuous,
        thm:mpu_complex_skew_unitary_normal_form,
        thm:mpug_kronecker_identity_factors}
    Let $u:\mathbb C^P\to\mathbb C^\ell\otimes\mathbb C^r$ and
    $v:\mathbb C^r\otimes\mathbb C^\ell\to\mathbb C^P$, and let $x,y$ be
    unitaries with
    \begin{align}
        \overline u & =(x\otimes y)u, & \overline v & =v(y^\dagger\otimes x^\dagger).
    \end{align}
    If $x=U^T\Lambda U$ and $y=V^T\Lambda'V$ with $U,V$ unitary, then
    $u'=(U\otimes V)u$ and $v'=v(V^\dagger\otimes U^\dagger)$ satisfy
    \begin{align}
        \overline{u'} & =(\Lambda\otimes\Lambda')u',                 &
        \overline{v'} & =v'(\Lambda'^\dagger\otimes\Lambda^\dagger).
    \end{align}
    Consequently, if $x,y$ are symmetric, then $u',v'$ can be chosen real; if
    $x,y$ are skew-symmetric, then $\ell,r$ are even and $u',v'$ can be chosen
    with
    \begin{align}
        \overline{u'} & =(\Sigma_\ell\otimes\Sigma_r)u', &
        \overline{v'} & =v'(\Sigma_r\otimes\Sigma_\ell),
    \end{align}
    where $\Sigma_x=\Id_{x/2}\otimes Y$.

    Conversely, let $\ell,r>0$, let $u'$ satisfy $u'u'^\dagger=\Id$ and
    $\overline{u'}=(\Lambda\otimes\Lambda')u'$ with $\Lambda,\Lambda'$
    unitary, and let $a,b$ be unitaries such that $(a\otimes b)u'$ satisfies
    the same equation. Then there is a phase $w$ with
    $\overline{wa}\,\Lambda=\Lambda\,wa$ and
    $\overline{w^{-1}b}\,\Lambda'=\Lambda'\,w^{-1}b$. For $\Lambda=\Lambda'=\Id$
    this gives real orthogonal $O_\ell=wa$, $O_r=w^{-1}b$ with
    $a\otimes b=O_\ell\otimes O_r$ and
    $b^\dagger\otimes a^\dagger=O_r^T\otimes O_\ell^T$. For
    $\Lambda=\Sigma_\ell$, $\Lambda'=\Sigma_r$ it gives unitaries
    $S_\ell=wa$, $S_r=w^{-1}b$ with $S_x^T\Sigma_xS_x=\Sigma_x$,
    $a\otimes b=S_\ell\otimes S_r$ and
    $b^\dagger\otimes a^\dagger=S_r^\dagger\otimes S_\ell^\dagger$.
    These are the computations in the proof of
    \cite[Proposition~V.2]{Cirac2017MPU}; the source prints $S_r^T\otimes
        S_\ell^T$ in the last relation, which is corrected in
    \cite{gap:mpu_conjugation_symplectic_uniqueness_transpose}.
    % Source lines: paper_v2.tex 1129--1163.
\end{lemma}

\begin{proof}
    \leanok
    \uses{lemma:conjclass-normalform-continuous,
        thm:mpu_complex_skew_unitary_normal_form,
        thm:mpug_kronecker_identity_factors}
    For a unitary $U$, $\overline UU^T=\Id$, hence
    $\overline U\,U^T\Lambda U=\Lambda U$ and
    $(U^T\Lambda U)^\dagger U^T=U^\dagger\Lambda^\dagger$. Conjugating $u'$
    and $v'$ and inserting the two equations gives the first claim. Symmetric
    unitaries are $S^TS$ with $S$ unitary, by
    Lemma~\ref{lemma:conjclass-normalform-continuous}; skew-symmetric
    unitaries have even dimension, since an invertible skew-symmetric matrix
    of odd size would have determinant zero, and are $U^T\Sigma U$ by
    Lemma~\ref{thm:mpu_complex_skew_unitary_normal_form}. Since
    $\Sigma^\dagger=-\Sigma$, the two signs cancel in
    $\Sigma_r^\dagger\otimes\Sigma_\ell^\dagger$.

    For the converse, conjugating $(a\otimes b)u'$ and cancelling $u'$ on the
    right gives
    $\overline a\Lambda\otimes\overline b\Lambda'=\Lambda a\otimes\Lambda' b$.
    Multiplying by the adjoint of the right side, a Kronecker product equals
    the identity, so $\overline a\Lambda=c\Lambda a$ and
    $\overline b\Lambda'=c^{-1}\Lambda'b$ for a scalar $c$, which has modulus
    one since both sides are unitary. A square root $w$ of $c$ has
    $\overline wc=w$, which gives the two intertwining relations. Finally
    $\overline S\Sigma=\Sigma S$ and $S^T\overline S=\Id$ give
    $S^T\Sigma S=\Sigma$.
\end{proof}

\begin{lemma}[The four-dimensional symplectic congruence]
    \label{lem:mpu_symplectic_four_congruence}
    \lean{Matrix.symplecticY, Matrix.symplecticOrthogonalR,
        Matrix.symplecticOrthogonalR_identities}
    \leanok
    Let $Y=\begin{pmatrix}0&1\\-1&0\end{pmatrix}$ and
    $R=\frac{1-i}{2}[\Id_4+i(Y\otimes Y)]$, with the two Kronecker factors in
    that order. Then
    \begin{align}
        R^\dagger R      & =\Id_4, \\
        R^T(Y\otimes Y)R & =\Id_4.
    \end{align}
    These are the identities of \cite[lines~1165--1185]{Cirac2017MPU}.
\end{lemma}

\begin{proof}
    \leanok
    Direct multiplication of the four entries of $Y$ and the corresponding
    entries of $Y\otimes Y$ gives both identities.
\end{proof}

\begin{lemma}[Reduction of the symplectic case to real matrices]
    \label{cor:sympl-to-orthogonal}
    \lean{Matrix.evenPairRegroupEquiv, Matrix.reindex_sourceSkewForm_kronecker,
        Matrix.symplecticOrthogonalR_aux,
        Matrix.exists_real_of_map_star_eq_sourceSkewForm_mul,
        Matrix.exists_real_of_map_star_eq_mul_sourceSkewForm}
    \leanok
    \uses{lem:mpu_symplectic_four_congruence}
    Let $\ell,r$ be even, and let $u',v'$ satisfy
    \begin{align}
        \overline{u'} & =(\Sigma_\ell\otimes\Sigma_r)u', &
        \overline{v'} & =v'(\Sigma_r\otimes\Sigma_\ell).
    \end{align}
    After regrouping coordinates one can write
    \begin{align}
        u' & =(\Id_{r\ell/4}\otimes R)\widetilde u, &
        v' & =\widetilde v(\Id_{r\ell/4}\otimes R),
    \end{align}
    where $\widetilde u$ and $\widetilde v$ are real,
    \begin{align}
        R & =\frac{1-i}{2}[\Id_4+i(Y\otimes Y)],
    \end{align}
    and $\widetilde u,\widetilde v$ are unitary between their coordinate
    spaces whenever $u',v'$ are. The Kronecker factors are ordered so that
    the two factors $Y$ of $\Sigma_\ell\otimes\Sigma_r$ act on the factor of
    $R$. This is \cite[Corollary~\texttt{cor:sympl-to-orthogonal}]{Cirac2017MPU},
    of which only the two displayed equations for $u',v'$ are used.
    % Source lines: paper_v2.tex 1165--1191.
\end{lemma}

\begin{proof}
    \leanok
    \uses{lem:mpu_symplectic_four_congruence}
    Regroup the coordinates so that
    $\Sigma_\ell\otimes\Sigma_r=\Id_{r\ell/4}\otimes(Y\otimes Y)$. The matrix
    $R$ is unitary and symmetric, and $R^T(Y\otimes Y)R=\Id_4$ gives
    $R^T(Y\otimes Y)=R^\dagger$ and $(Y\otimes Y)R^T=R^\dagger$. Set
    $\widetilde u=(\Id\otimes R^\dagger)u'$ and
    $\widetilde v=v'(\Id\otimes R^\dagger)$. Then
    \begin{align}
        \overline{\widetilde u} & =(\Id\otimes R^T)(\Id\otimes Y\otimes Y)u'=\widetilde u, \\
        \overline{\widetilde v} & =v'(\Id\otimes Y\otimes Y)(\Id\otimes R^T)=\widetilde v,
    \end{align}
    and $\widetilde u,\widetilde v$ are unitary whenever $u',v'$ are.
\end{proof}

\begin{lemma}[An involution converts a combined fixed equation]
    \label{thm:mpu_involution_combined_fixed_equation}
    \lean{Matrix.swapCombined_fixed_iff_leftMul}
    \leanok
    Let $Q,X$ be complex matrices indexed by the same finite set, and let
    $F$ be a map on that matrix space. If $Q^2=I$ and
    $F(QX)=F(X)Q$, then
    \begin{align}
        QF(X)Q=X
         & \quad\Longleftrightarrow\quad F(QX)=QX.
        \label{eq:mpu_involution_combined_fixed}
    \end{align}
    This is the matrix argument of \cite[lines~1295--1322]{Cirac2017MPU}.
\end{lemma}

\begin{proof}\leanok
    Replacing $F(QX)$ by $F(X)Q$ and multiplying on the left by the
    involution $Q$ gives both implications in
    (\ref{eq:mpu_involution_combined_fixed}).
\end{proof}
